\section{Spectral solution of the Lyapunov equation and transient protocol integrals}
\quad This appendix collects two technical steps used in the main text: the spectral solution of the Lyapunov equation and the mode expansion of the time-dependent convolution integrals in Sec. \ref{sec:time-dependent}.
\subsection*{Spectral solution of the Lyapunov equation}
Consider the Lyapunov equation for the unknown matrix ${\bf K}$ \cite{Behr2019SylvesterLyapunov}
\begin{equation}
{\bf Q}{\bf K}+{\bf K}{\bf Q}^\intercal={\bf S},
\label{lyap}
\end{equation}
where ${\bf Q}$ and ${\bf S}$ are given matrices and ${\bf S}$ is symmetric.
Assume that ${\bf Q}$ is diagonalizable:
\begin{equation}
{\bf Q}={\bf P}\boldsymbol{\Lambda}{\bf P}^{-1},\quad
{\bf P}=({\vec u_1},{\vec u_2},\cdots,{\vec u_N}),\quad
\boldsymbol{\Lambda}=\hbox{diag}(\lambda_1,\lambda_2,\cdots,\lambda_N),
\end{equation}
with ${\bf Q}{\vec u_i}=\lambda_i{\vec u_i}$. Let the columns of
${\bf P}^{-\intercal}$ be the left eigenvectors,
\begin{equation}
{\bf P}^{-\intercal}=({\vec v_1},{\vec v_2},\cdots,{\vec v_N}),
\end{equation}
so that
\begin{equation}
{\vec v_i}^\intercal{\vec u_j}=\delta_{ij},\qquad
\sum_{j=1}^N {\vec u_j}{\vec v_j}^\intercal={\bf I}.
\end{equation}
Defining ${\tilde{\bf K}}={\bf P}^{-1}{\bf K}{\bf P}^{-\intercal}$, Eq. (\ref{lyap}) becomes
\begin{eqnarray}
\boldsymbol{\Lambda}{\tilde{\bf K}}+{\tilde{\bf K}}\boldsymbol{\Lambda}
&=&{\bf P}^{-1}{\bf S}{\bf P}^{-\intercal},\\
{\tilde{\bf K}}_{ij}
&=&\frac{({\bf P}^{-1}{\bf S}{\bf P}^{-\intercal})_{ij}}{\lambda_i+\lambda_j}.
\label{Ktil}
\end{eqnarray}
Hence
\begin{equation}
{\bf K}={\bf P}{\tilde{\bf K}}{\bf P}^\intercal.
\end{equation}
For the stable fixed points considered here, $\mathrm{Re}\,\lambda_i<0$, so $\lambda_i+\lambda_j\neq 0$ in Eq. (\ref{Ktil}). Using the biorthogonal decomposition
\begin{equation}
{\bf Q}=\sum_{m=1}^N \lambda_m {\vec u_m}{\vec v_m}^\intercal,
\end{equation}
the same solution can be written more compactly as
\begin{equation}
{\bf K}=\sum_{m,\ell=1}^N
\frac{({\vec v_m}^\intercal{\bf S}{\vec v_\ell})\,{\vec u_m}{\vec u_\ell}^\intercal}
{\lambda_m+\lambda_\ell}.
\end{equation}
\subsection*{Mode expansion for the time-dependent protocol integrals}
The time-dependent formulas in Sec. \ref{sec:time-dependent} follow from the decomposition of $\mathbf{Q}$. Since
\begin{equation}
e^{t{\bf Q}}=\sum_{k=1}^N e^{\lambda_k t}{\vec u_k}{\vec v_k}^\intercal,
\qquad
{\bf Q}^{-1}=\sum_{k=1}^N \frac{{\vec u_k}{\vec v_k}^\intercal}{\lambda_k},
\end{equation}
Eq. (\ref{xi(t)drlinear}) becomes
\begin{equation}
\langle x_i(t)\rangle_{r+\delta r_\alpha(t)}=X_i+(a_\alpha-X_\alpha)\bigg[
\sum_{k=1}^N \frac{u_{ik}v_{\alpha k}}{\lambda_k}\int_0^t dt'
e^{(t-t')\lambda_k}\delta{\dot r}_\alpha(t')
-Q^{-1}_{i\alpha}\delta r_\alpha(t)\bigg],
\label{xi(t)drlinearb}
\end{equation}
where $u_{ik}\equiv({\vec u_k})_i$ and $v_{\alpha k}\equiv({\vec v_k})_\alpha$.
The time-dependent relaxation response is then
\begin{equation}
      \langle x_i(t)\rangle=\begin{cases}  X_i+ \frac{b}{2} (a_\alpha-X_\alpha)\bigg[\sum_{k=1}^N\frac{u_{ik}v_{\alpha k}}{\lambda_k(1+(\frac{\lambda_k\tau}{\pi})^2)}\left(e^{\lambda_kt}-\cos\frac{\pi t}{\tau}-\frac{\lambda_k\tau}{\pi}\sin\frac{\pi t}{\tau}  \right)-&Q^{-1}_{i\alpha} (1-\cos\frac{\pi t}{\tau})\bigg] \\
     & 0\leq t\leq\tau\\
     X_i-b (a_\alpha-X_\alpha)Q^{-1}_{i\alpha}+\frac{b}{2} (a_\alpha-X_\alpha)\sum_{k=1}^N\frac{u_{ik}v_{\alpha k}[ e^{\lambda_kt}+e^{\lambda_k(t-\tau)}]}{\lambda_k(1+(\frac{\lambda_k\tau}{\pi})^2)} &t>\tau
     \end{cases}.
\end{equation}
For the force response,
\begin{equation}
\langle { F}_i(t)\rangle = \begin{cases}  (a_\alpha-X_\alpha)\sum_{k=1}^N u_{ik}v_{\alpha k}\int_0^tdt' \delta {\dot r}_\alpha(t')e^{(t-t')\lambda_k}, &0\leq t\leq\tau\\
(a_\alpha-X_\alpha)\sum_{k=1}^N u_{ik}v_{\alpha k}\int_0^\tau dt' \delta {\dot r}_\alpha(t')e^{(t-t')\lambda_k}, & t>\tau
\end{cases}
\end{equation}
\begin{equation}
\langle { F}_i(t)\rangle= \begin{cases} \frac{b}{2}(a_\alpha-X_\alpha)\sum_{k=1}^N \frac{u_{ik}v_{\alpha k}}{1+\left(\frac{\lambda_k\tau}{\pi}\right)^2}\left[ e^{\lambda_k t} - \cos\left( \frac{\pi t}{\tau} \right) - \frac{\lambda_k \tau}{\pi}\sin\left( \frac{\pi t}{\tau} \right) \right], &0\leq t\leq\tau\\
\frac{b}{2}(a_\alpha-X_\alpha)\sum_{k=1}^N \frac{u_{ik}v_{\alpha k}}{1+\left(\frac{\lambda_k\tau}{\pi}\right)^2} \left[ e^{\lambda_k t}+e^{\lambda_k(t-\tau)} \right], & t>\tau
\end{cases}.
\end{equation}
Likewise, for the covariance response,
\begin{eqnarray}
&&\int_0^t dt' \left(
e^{(t-t'){\bf Q}}
\frac{\partial {\bf K}_0}{\partial Q_{\alpha\alpha}}
e^{(t-t'){\bf Q}^\intercal}
\right)_{ij}\delta \dot{r}_\alpha(t') \nonumber\\
&=&\sum_{m,\ell=1}^N
u_{im}u_{j\ell}\left(
{\vec v_m}^\intercal
\frac{\partial {\bf K}_0}{\partial Q_{\alpha\alpha}}
{\vec v_\ell}
\right)I_{m\ell}(t),
\end{eqnarray}
with
\begin{equation}
I_{m\ell}(t)\equiv
\int_0^t dt' \, e^{(t-t')(\lambda_m+\lambda_\ell)}\delta \dot r_\alpha(t')
=\frac{b\pi}{2\tau}\int_0^{\min(t,\tau)} dt' \, e^{(t-t')(\lambda_m+\lambda_\ell)}\sin\left(\frac{\pi t'}{\tau}\right),
\label{Iml}
\end{equation}
which evaluates to
\begin{equation}
I_{m\ell}(t)=
\begin{cases}
\displaystyle
\frac{b}{2\left[1+\left(\frac{(\lambda_m+\lambda_\ell)\tau}{\pi}\right)^2\right]}
\left[
e^{(\lambda_m+\lambda_\ell)t}
-\cos\left(\frac{\pi t}{\tau}\right)
-\frac{(\lambda_m+\lambda_\ell)\tau}{\pi}\sin\left(\frac{\pi t}{\tau}\right)
\right], & 0\le t\le \tau, \\[1.2ex]
\displaystyle
\frac{b}{2\left[1+\left(\frac{(\lambda_m+\lambda_\ell)\tau}{\pi}\right)^2\right]}
\left[
e^{(\lambda_m+\lambda_\ell)t}
+e^{(\lambda_m+\lambda_\ell)(t-\tau)}
\right], & t>\tau .
\end{cases}
\end{equation}
The edge-perturbation formulas in Sec. \ref{sec:time-dependent} follow from the same expansion, with only the prefactor and the covariance derivative replaced according to the perturbation under consideration.

The corresponding time-dependent response of each node state variable is
\begin{equation}
      \langle x_i(t)\rangle=\begin{cases}  X_i+ \frac{b}{2} (X_\beta-X_\alpha)\bigg[\sum_{k=1}^N\frac{u_{ik}v_{\alpha k}}{\lambda_k(1+(\frac{\lambda_k\tau}{\pi})^2)}\left(e^{\lambda_kt}-\cos\frac{\pi t}{\tau}-\frac{\lambda_k\tau}{\pi}\sin\frac{\pi t}{\tau}  \right)&-Q^{-1}_{i\alpha} \left(1-\cos\frac{\pi t}{\tau}\right)\bigg] \\
     & 0\leq t\leq\tau\\
     X_i-b (X_\beta-X_\alpha)Q^{-1}_{i\alpha}+\frac{b}{2} (X_\beta-X_\alpha)\sum_{k=1}^N\frac{u_{ik}v_{\alpha k}[ e^{\lambda_kt}+e^{\lambda_k(t-\tau)}]}{\lambda_k(1+(\frac{\lambda_k\tau}{\pi})^2)} &t>\tau
     \end{cases}.
\end{equation}

\begin{equation}
\langle { F}_i(t)\rangle= \begin{cases} \frac{b}{2}(X_\beta-X_\alpha)\sum_{k=1}^N \frac{u_{ik}v_{\alpha k}}{1+\left(\frac{\lambda_k\tau}{\pi}\right)^2}\left[ e^{\lambda_k t} - \cos\left( \frac{\pi t}{\tau} \right) - \frac{\lambda_k \tau}{\pi}\sin\left( \frac{\pi t}{\tau} \right) \right], &0\leq t\leq\tau\\
\frac{b}{2}(X_\beta-X_\alpha)\sum_{k=1}^N \frac{u_{ik}v_{\alpha k}}{1+\left(\frac{\lambda_k\tau}{\pi}\right)^2} \left[ e^{\lambda_k t}+e^{\lambda_k(t-\tau)} \right], & t>\tau
\end{cases}.
\end{equation}

The time-dependent response function of the covariance is
\begin{eqnarray}
    {\bf K}_{ij}(t,t) &=& ({\bf K_0})_{ij} + \chi_{ij\alpha\beta}^{ss} \delta W_{\alpha\beta}(t)+\int_0^t dt' \left[ e^{(t-t'){\bf Q}}\left(-\frac{\partial {\bf K}_0}{\partial Q_{\alpha\alpha}} +\frac{\partial {\bf K}_0}{\partial Q_{\alpha\beta}}\right) e^{(t-t'){\bf Q}^\intercal} \right]_{ij} \delta \dot{W}_{\alpha\beta}(t')\nonumber\\
    &=&({\bf K_0})_{ij}+\left[- \frac{\partial {\bf K_0}}{\partial Q_{\alpha\alpha}}+\frac{\partial {\bf K_0}}{\partial Q_{\alpha\beta}} \right]_{ij}\delta W_{\alpha\beta}(t)\nonumber\\
    &+& \sum_{m,\ell} u_{im}u_{j\ell} \left[ \vec{v}^\intercal_m \left( \frac{\partial {\bf K_0}}{\partial Q_{\alpha\alpha}}-\frac{\partial {\bf K_0}}{\partial Q_{\alpha\beta}} \right) \vec{v}_\ell \right] I_{m\ell}(t).
\end{eqnarray}
